% \begin{savequote}
% \includegraphics[width=5cm]{./images/motal-aass/00800/neuro2.jpg}
% 
% Bild: Motal
% \end{savequote}
\def\ChapterCode{BAU}
\chapter
[\CO{[BAU]} Basisuntersuchungen]
{Basisuntersuchungen}
\label{chp:BAU}

\ChapterIntro
{%
\noindent Dieses Kapitel behandelt die Basisuntersuchungen.
}
{%
\begin{description}
  \item[Maintainer:] Sebastian Gabriel
  \item[Autoren:] Diverse
  \item[Reviewer:] --
  \item[Version:] {\DocVersion} ({\DocVersionInfo})
  \item[SHA1:] {\DocGitCommit}
\end{description}

{\TxTeilkapitelAASS}
}





\Dictum
[Mephistopheles. \\
In: J. W. von Goethe: Faust. Der Tragödie zweiter Teil]
{Zwar ist es leicht,\\
doch ist das Leichte schwer.}

\clearpage % TEMP temp clearpage
% \TextSlide{\TxQuerverweise}
% {~
% 
% \noindent\includegraphics[angle=270,width=\linewidth]{./images/leitsystem/Leitschema-PAM_BAU.pdf}
% }{\begin{itemize}
%   \item \EE{Medizinprodukte allgemein}: \CO{[MP1]} (\ref{chp:MP1})
%   \item \EE{Untersuchungen}:                 \CO{[BAU]} (\ref{chp:BAU})
%   \item \EE{Sanitätshilfemaßnahmen}:         \CO{[SAN]} (\ref{chp:SAN})
%   \item \EE{Anamnese und Patientengespräch}: \CO{[AUP]} (\ref{chp:AUP})
% \end{itemize}}{}{}{}{}
% 
% 
% 
% 
% 





\section{Vitalwerte}
\label{topic:vitalwerte}

\TextSlide{\TxBeschreibung}{
Das Erheben von Vitalwerten ist eine Routinetätigkeit und
sollte bei jedem Notfallpatienten erfolgen. Es gibt nur wenig
Situationen bei denen man darauf verzichten kann (zum
Beispiel tobende Patienten). Zu den Vitalwerten zählt man:

\begin{itemize}
  \item die \EE{Herzfrequenz} {\ldots}
  \FullPrettyRef{topic:pulsmessung}, \FullPrettyRef{topic:pulsoxymetrie}
  \item den \EE{Blutdruck} {\ldots} \FullPrettyRef{topic:rr-messung}
  \item die \II{Atemfrequenz}: Bei der Messung der
  Atemfrequenz wird die Anzahl der Atemzüge pro Minute
  gezählt.
\end{itemize}
}{
\begin{itemize}
  \item Herzfrequenz
  \item Blutdruck
  \item Atemfrequenz
\end{itemize}
}{}{}{}



\subsection{Ermitteln der Herzfrequenz mittels Messung der
Pulsfrequenz}
\label{topic:pulsmessung}


\TextSlide
{{\TxBeschreibung} und Messstellen}
{%
\DefinitionInPara{Als \T{Herzfrequenz} (\T{HF})
bezeichnet man die Anzahl der Schläge bzw. Aktionen des
Herzens pro Minute. Als \T{Pulsfrequenz} bezeichnet man die
Anzahl der Pulsschläge pro Minute.}
%
Die Herzfrequenz entspricht normalerweise der
\T{Pulsfrequenz}, welche mittels Fühlen des Pulses 
ermittelt wird.\footnote{Nicht
immer hat eine Herzaktion einen Pulsschlag zur Folge, man
spricht dann vom \I{Pulsdefizit}. Außerdem kann man
zwischen elektrischen und mechanischen Herzaktionen
unterscheiden (vgl. \FullPrettyRef{topic:herz}).}
Dieser wird idealerweise an der
\EE{Radialarterie} (A. \E{radialis}) bzw. bei nicht
ansprechbaren oder schockierten Patienten an der
\EE{Halsschlagader} (A. \E{carotis}) ertastet. Bei
Neugeborenen und Säuglingen wird an der
\EE{Oberarmarterie} (A. \E{brachialis}) gefühlt. Gemessen
wird immer mit zwei bis drei Fingern, nie jedoch mit dem
Daumen, da hier oft irrtümlich der \so{eigene} Puls gefühlt
wird, und nicht der des Patienten. Wenn an der
Halsschlagader (A. \E{carotis}) die Herzfrequenz ermittelt
wird, darf nie auf beiden Seiten gleichzeitig gemessen
werden, da sonst die Gefahr besteht,
dass es zu einer Unterversorgung des Gehirn durch Blut kommen
kann.
}
{%&
\begin{itemize}
  \item Herzfrequenz, Pulsfrequenz
  \item {Radialarterie} (A. \E{radialis})
  \item {Halsschlagader} (A. \E{carotis})
  \item {Oberarmarterie} (A. \E{brachialis}; Neugeborene, Säuglinge)
\end{itemize}
}
{}
{}
{}

\begin{PictureSeriesWide}{}{Ermittlung der Pulsfrequenz \SourcePicture{ASBÖ/gjh}}
\PictureSeriesRowWide{3}
{./images/ASBOE-Grassl-gjh-1/radialis_puls-00441pt}
{Tasten des Radialispuls mit mind. 2 Fingern}
{./images/ASBOE-Grassl-gjh-1/rate_puls-00441pt}
{Auszählen der Pulsfrequenz (Anm.: Entgegen der Darstellung ist eine Armbanduhr im Dienst aus hygienischen Gründen \EE{nicht} am Handgelenk zu tragen!)}
{./images/ASBOE-Grassl-gjh-1/puls_oximeter-00441pt}
{Die Pulsoxymetrie kann oft bei der Ermittlung der Pulsfrequenz behilflich sein {\ldots} jedoch gibt es hier Tücken}
{}
{}
\end{PictureSeriesWide}







\TextSlide{Durchführung}{
Es wird an der Arterie (s.\,o.) der Puls mit
zwei bis drei Fingern ertastet. Nun werden 15 Sekunden lang
die Schläge gezählt. Anschließend wird das Ergebnis \EE{mit
vier multipliziert}, man erhält die Herzfrequenz (HF) in
Pulsschläge \EE{pro Minute} (\PerMin).
Der Normalwert liegt beim Erwachsenen zwischen 60 und 100
Schlägen pro Minute. Neben der Frequenz wird auch die
\EE{Regelmäßigkeit} beurteilt, das heißt ob der Puls
\E{rhythmisch} oder \E{arrhythmisch} ist.

}{%
\begin{itemize}
  \item 15\Sec*
  \item \Range{2}{3} Finger
  \item Mit 4 multiplizieren 
\end{itemize}
}{}{}{}

\TextSlide{Befunde}{
\begin{Findings}
\item[\SymbolFindingNormal] \FindingSpeech{Der Radialispuls ist gut tastbar, rhythmisch und normfrequent} (\FindingSpeech{und hat eine Frequenz von 80\PerMin*}) [\Range{60}{100}\PerMin*]
\item[\SymbolFindingVariant] \FindingSpeech{Der Carotispuls \ldots}
\end{Findings}

}{%
\begin{Findings}
\item[\SymbolFindingNormal] \FindingSpeech{Radialispuls gut tastbar, rhythmisch, normfrequent} [\Range{60}{100}\PerMin*]
\item[\SymbolFindingVariant] \FindingSpeech{Der Carotispuls \ldots}
\end{Findings}
}{}{}{}


\subsection{Messung des Blutdrucks -- RR}
\label{topic:rr-messung}
\label{topic:blutdruckmessung}

\TextSlide{\TxBeschreibung}{
Der Blutdruck ist jener Druck, welcher in unseren Blutgefäßen
herrscht. Wenn wir vom \emph{\enquote{Blutdruck}} sprechen beziehen wir
uns auf den Blutdruck in den
Arterien\footnote{\emph{Arterie:} Vom Herzen wegführendes
Blutgefäß.}. Er gehört in Kombination mit der Herzfrequenz
zu den grundlegenden Vitalwerten.

\E{Der Druck in der Arterie schwankt in Abhängigkeit zum
Herzauswurf}: Wenn das Herz Blut auswirft
(\E{Auswurfphase}, \T{Systole}) erreicht er sein Maximum,
und sinkt danach in der \E{Füllungsphase} des Herzens
(\T{Diastole}) wieder ab, bis das Herz neuerlich Blut auswirft.
Dementsprechend können zwei Werte erhoben werden: Der

\begin{itemize}
  \item \T{Systolische Blutdruck} (Spitzendruck), und der
  \item \T{Diastolische Blutdruck} (Minimaldruck).
\end{itemize}
Im Normalfall
werden beide Werte ermittelt.
% Der \EE{systolische Wert} (\enquote{obere Wert}) ist jener Druck
% im Blutgefäßsystem, der  während der \E{Auswurfphase} des
% Herzens, in der sich die Herzkammern anspannen,
% zusammenziehen und Blut auswerfen, herrscht. Der
% \EE{diastolische Wert} (\enquote{untere Wert}) ist jener Druck im
% Blutgefäßsystem, der während der \E{Entspannungsphase} des
% Herzens, in der sich die Herzkammern mit Blut füllen,
% herrscht.
Die \EE{Einheit}, in der der Blutdruck angegeben wird, ist
\T{Millimeter Quecksilbersäule}
(\T{mm\,Hg}).\index{Millimeter Quecksilbersäule}\index{mm\,Hg}\footnote{\ce{Hg}
ist das chemische Symbol für Quecksilber.}
}{
\begin{itemize}
  \item Druck in den Arterien
  \begin{itemize}
    \item Systolischer Wert (Auswurfphase des Herzens)
    \item Diastolischer Wert (Entspannungsphase des Herzens)
  \end{itemize}
  \item Einheit: Millimeter Quecksilbersäule (\MmHg)
\end{itemize}
}{}{}{}

\TextSlide{Arten der Blutdruckmessung}{%
Der Blutdruck kann auf verschiedene Arten gemessen werden.
Üblich ist die unblutige, \E{indirekte Methode}, mittels
einer an der Extremität angelegten
Druckmanschette\ZFootnote{Mittels einer Druckmanschette wird
gemessen, bei welchem Druck ein mit einem Stethoskop
hörbarer, durch Pulstasten fühlbarer oder durch sonstige
technische Hilfsmittel messbarer Blutfluss stattfindet.
Dies erlaubt Rückschlüsse auf den Druck im Blutgefäß.}. Dazu
gibt es unterschiedliche Messgeräte, grundsätzlich
unterscheidet man zwischen elektronischen und rein
mechanischen Messgeräten.

\EE{Elektronische Messgeräte}  werden oft von Patienten zur
Selbstkontrolle des Blutdrucks verwendet. Diese Geräte sind
jedoch vergleichsweise fehleranfällig und sind für einen
professionellen Einsatz nicht geeignet.

Für den professionellen Einsatz gibt es spezielle
elektronische Blutdruckmessgeräte, welche zum Beispiel in
Überwachungsmonitoren zum Einsatz kommen. Allerdings sind
auch diese Geräte mitunter fehleranfällig, wenn die
Messbedingungen nicht optimal (z.\,B. bei Erschütterungen)
sind.

% Der Blutdruck kann mittels elektronischer
% Blutdruckmessgeräte ermittelt werden. Diese sind aber im
% Rettungsdienst nicht in Verwendung, da es immer wieder zu
% Gerätefehler kommen kann. Außerdem können die Batterien leer
% werden. Solche Geräte haben meistens die Patienten für die
% tägliche Überwachung ihres Blutdrucks. Leider bedienen aber
% viele die Geräte falsch und fürchten sich dann
% lebensbedrohlich krank zu sein.

In der Praxis wird die Blutdruckmessung mit \EE{mechanischen
Messgeräten} durchgeführt. Dabei wird eine aufblasbare
Manschette mit angeschlossenem Dosenmanometer und Pumpballen
mit einem Ablassventil benötigt. Im Idealfall wird auch noch
ein Stethoskop verwendet.

\begin{center}
     
%   \includegraphics[width=4cm]{./images/noauth/GER-rr-elektrisch-oberarm.png}
%   
%    \includegraphics[width=4cm]{./images/noauth/GER-rr-elektrisch-unterarm.jpg}
%    
%    \includegraphics[width=4cm]{./images/noauth/GER-rr-messgeraet.jpg}
\TODO{GABS}{2010-04-27}{Blutdruckmesser: Bilder aus Serie
Pallinger}{}
\end{center}
}{
\begin{itemize}
  \item Unblutige, indirekte Messung
  \begin{itemize}
    \item Mechanische Messung
    \item Elektronische Messung
  \end{itemize}
  \item Invasive Messung
\end{itemize}
}
{}
{}
{}


\subsubsection{Manuelle Blutdruckmessgeräte}
\Text{Manuelles Blutdruckmessgerät}{%
Ein manuelles Blutdruckmessgerät besteht aus einer
Manschette, einem Schlauch, einem Auslassventil, einem Pumpballen und einer
Druckanzeige. 
Es gibt Manschetten in unterschiedlichen Größen.
}
{}
{}
{}
{}

\Layout{37}{Abbildung}
{}
{}
{./images/pallinger-christoph-aass/Blutdruck_32768_export_01024.jpg}
{Blutdruckmanschetten mit Hak- und Klettverschluss, sowie
einem Doppel- und einem Einzelkopfstethoskop.}
{Ch. Pallinger}

\subsubsection{Exkurs: Stethoskop}
\label{topic:stethoskop}
\index{Stethoskop}

\Text{\TxBeschreibung}
{
\begin{description}
  \item[Allgemein] Ein Stethoskop ist eine Hörvorrichtung. Es besteht
patientenseitig aus einem Kopf mit einer Membran oder einem
Trichter, einem Schlauch und untersucherseitig aus einem
Bügel mit zwei Oliven (\enquote{Ohrenstöpsel}).
Die Oliven sollen \E{nach vorne} gerichtet sein, da auch der
äußere Gehörgang nach vorne verläuft.
\item[Funktionen] Manche Typen verfügen über einen
zweiseitigen Kopf.
Meistens haben sie sowohl eine Membran, als auch einen
Schalltrichter.~\ZFootnote{Der Unterschied besteht in einer
unterschiedlichen Verstärkung von hohen oder tiefen
Frequenzen, was für den fortgeschrittenen Anwender beim
Abhören von Herz oder Lunge eine Rolle spielt.
Teurere Modelle verfügen mitunter über eine
Dual-Frequency-Membran, mit welcher man durch den
Auflegedruck wahlweise hohe oder tiefe Frequenzen verstärken
kann.} Durch Drehen am Kopf kann man die gewünschte Seite
einstellen.
\item[Funktionsprüfung] Vor der Verwendung muss mittels Beklopfen der Membran oder
des Schalltrichters geprüft werden, ob der Kopf des
Stethoskops auf die gewünschte Seite eingestellt ist und ob
die Schallleitung funktioniert.
\item[Hygienische Wiederaufbereitung] Das Stethoskop ist ein
häufiges Vehikel zur Keimübertragung. Nach jedem Patienten soll eine 
Wischdesinfektion mit Flächendesinfektionsmittel
vorgenommen werden.
\end{description}

}{}{}{}{}



\subsubsection{Messung nach Riva-Rocci}

\Text{{\TxBeschreibung}}{%
\DefinitionInPara{Die \T{Blutdruckmessung nach
Riva-Rocci}\ZFootnote{Benannt nach dem Erfinder Scipione
Riva-Rocci.} ist eine indirekte, unblutige, mechanische 
Messung mittels einer
Druckmanschette.}
Umgangssprachlich wird die Blutdruckmessung nach Riva-Rocci
als \emph{\enquote{RR-Messung}}, bzw. der ermittelte Wert
als \emph{\enquote{RR}} abgekürzt.

Bei der RR-Messung gibt es zwei Möglichkeiten zu
einem Wert zu kommen. Die genauere ist die
\EE{auskultatorische Methode} (durch Hören von
Strömungsgeräuschen). Die zweite Möglichkeit ist die
\EE{palpatorische Methode} (durch Tasten des Pulses). 
}

\Fazit{Der diastolische Wert kann nur
bei der auskultatorischen Methode ermittelt werden!} 


\TextSlide{\T{Auskultatorische Methode} -- Vorgehen}
{%
Die Manschette sollte faltenfrei und nicht über Kleidung, in
der Mitte des Oberarms angelegt werden und sich, wenn
möglich, auf \E{Herzhöhe} befinden\ZFootnote{Wenn die
Manschette nicht auf Herzhöhe liegt, wird ein zu hoher oder
zu niedriger Druck gemessen.} Die Mitte der Luftkammer soll
über der Oberarmarterie zu liegen kommen. Manche Manschetten
haben einen Markierungspfeil für die Oberarmarterie, der
Pfeil soll dann oberhalb dieser zu liegen kommen. Man achte
auf die Lage der Innen- bzw. Außenseite!

Nun wird der Puls an der Radialarterie (A. radialis)
getastet. Die Manschette wird solange aufgepumpt, bis der
Puls nicht mehr getastet werden kann. Dann werden noch
\VonBis{20}{30}\MmHg* weiter aufgepumpt. Das Stethoskop wird
in der Ellenbeuge auf die Oberarmarterie (A. brachialis,
Verlauf eher zur Körpermitte) angelegt. Die Luft wird
langsam aus der Manschette abgelassen.

Der Wert, bei dem man erstmalig ein Pochen \Z{(Korotkowsche
Geräusche)} hört, gilt als \EE{systolischer Wert}. Die Luft
wird weiter abgelassen. Der Wert, bei dem das Pochen
verschwindet, ist der \EE{diastolische Wert}\footnote{Es
kann vorkommen, dass das Pochen nie verschwindet, sondern
nur plötzlich leiser wird. Dann gilt der Druck, bei dem das
Pochen \emph{deutlich} leiser wird, als diastolischer
Wert.}. \cite{FamPropAeGruFert2012}

Die
Werte werden auf die nächste Fünferstelle abgerundet
(z.\,B.:
\RRmmHg{122}{82} wird abgerundet auf \RRmmHg{120}{80}).
Dokumentiert wird der Blutdruck mit dem Zusatz \E{RR} (z.\,B.:
RR~\RRmmHg{120}{80}). Der Normalwert beim Erwachsenen liegt
in etwa bei RR~\RRmmHg{120}{80}
(\prettyref{topic:blutdruck}).
}{
\begin{compactenum}
  \item Manschette faltenfrei auf Herzhöhe anlegen
  \item Luftkissen  mittig über Oberarmarterie
  (Markierungspfeil!)
  \item Radialarterie tasten
  \item Manschette aufpumpen bis Radialispuls nicht mehr
  tastbar ist +~\VonBis{20}{30}\,\MmHg
  \item Stethoskop über Oberarmarterie anlegen
  \item Druck langsam ablassen
  \item Erstes Pochen: systolischer Wert
  \item Letztes Pochen: diastolischer Wert
  \item Werte auf nächste 5er-Stelle abrunden
\end{compactenum}

}{}{}{}


\begin{PictureSeriesWide}{}{Blutdruckmessung \SourcePicture{ASBÖ/gjh}}
\PictureSeriesRowWide{3}
{./images/ASBOE-Grassl-gjh-1/rekap-00441pt}
{Die Nagelbettprobe (Rekap-Zeit) dient einer schnellen Einschätzung der Kreislaufsituation. Normal ist eine Rekap-Zeit von unter 2 Sekunden.}
{./images/ASBOE-Grassl-gjh-1/rr_messung-00441pt}
{Blutdruckmessung}
{./images/ASBOE-Grassl-gjh-1/CRW_0784-00441pt}
{Blutdruckmessung bei einem Kind}
{}
{}
\end{PictureSeriesWide}

\begin{figure}[H]
    \begin{center}
    \includegraphics[width=\linewidth]{./images/hirtler-aass/Riva-Rocci_edited.pdf}
   
    \Captionof{figure}{Prinzip
    der auskultatorischen RR-Messung: In Abhängigkeit vom
    außen anliegenden Manschettendruck kommt es zu einem
    unterschiedlich starken Blutfluss und dadurch
    entstehenden, hörbaren, Strömungsgeräuschen.}
    \end{center}
\end{figure}


\TextSlide{\T{Palpatorische Methode} -- Durchführung}
{%
Die Manschette sollte faltenfrei und nicht über Kleidung, in
der Mitte des Oberarms angelegt werden und sich, wenn
möglich, auf \E{Herzhöhe} befinden\ZFootnote{Wenn die
Manschette nicht auf Herzhöhe liegt wird ein zu hoher oder
zu niedriger Druck gemessen.}. Nun wird der Puls an der
Radialarterie (A. radialis) ertastet. Die Manschette wird
solange aufgepumpt, bis der Puls nicht mehr getastet werden
kann, dann wird der Druck um weitere 20\MmHg* erhöht. Die
Finger bleiben auf der Radialarterie (A.
radialis). Die Luft wird langsam aus der Manschette
abgelassen. Der Druck, bei dem der Puls wieder fühlbar ist,
ist der \EE{systolische} Wert.

\EE{Der diastolische Wert kann mit dieser Methode nicht
ermittelt werden.} In der Notfallmedizin ist der systolische
Wert jedoch ohnehin in der Regel der wichtigere.

\bigskip
}{
\begin{compactenum}
  \item Manschette faltenfrei auf Herzhöhe anlegen
  \item Luftkissen  mittig über Oberarmarterie
  (Markierungspfeil!) beachten
  \item Radialarterie tasten
  \item Manschette aufpumpen bis Radialispuls nicht mehr tastbar ist
  +~\VonBis{20}{30}\,\MmHg
  \item Druck langsam ablassen
  \item Radialispuls wieder fühlbar: systolischer Wert
  \item \E{Kein} \E{diastolischer} Wert!
\end{compactenum}
}{}{}{}



\Layout{60}{Was ist besonders zu beachten?}
{%
Die Blutdruckmanschette soll nach Möglichkeit \EE{\E{nicht}
auf der Seite
\ldots}
\begin{inparaenum}
  \item eines \EE{Shuntarmes}\footnote{Ein
  (Dialyse-)\EE{Shunt} ist ein künstlicher Kurzschluss
  zwischen einer Armarterie und einer Armvene. Dadurch wird
  die Vene erweitert und sie pulsiert. Man kann viel mehr
  Durchfluss für die Dialyse erreichen. Die Nahtstelle
  zwischen Vene und Arterie ist jedoch sehr sensibel. Wenn
  gestaut wird, kann sie aufplatzen $\rightarrow$ Gefahr
  einer massiven Shuntblutung.} bei Dialysepatienten \item
  eines Armödems, z.\,B. nach einer Brustamputation, oder
  \item einer Halbseitenschwäche oder -lähmung
\end{inparaenum}
angelegt werden.

\Cave{Shunt-Arme sind tabu für die Blutdruckmessung!}
}{
\begin{itemize}
  \item \EE{Nicht} auf Seite eines/r \dots{}
  \begin{enumerate}
    \item \EE{Shuntarmes}
    \item Armödems
    \item Halbseitenschwäche/-lähmung
  \end{enumerate}
\end{itemize}

}{./images/sign_cardive.jpg}{}{Gabriel}


\TextSlide{Befunde}{
\begin{Findings}
\item[\SymbolFindingNormal] \FindingSpeech{Blutdruck (RR) \RRmmHg{120}{80}}  [\Range{\RR{100}{60}}{\RR{140}{90}}\MmHg*]
\item[\SymbolFindingVariant] \FindingSpeech{Blutdruck (RR) palpatorisch 120\MmHg*}  [\Range{100}{140}\MmHg*]
\end{Findings}

}{%
\begin{Findings}
\item[\SymbolFindingNormal] \FindingSpeech{RR \RRmmHg{120}{80}}  [\Range{\RR{100}{60}}{\RR{140}{90}}\MmHg*]
\item[\SymbolFindingVariant] \FindingSpeech{RR palp. 120\MmHg*}  [\Range{100}{140}\MmHg*]
\end{Findings}
}{}{}{}




\Z{
\Text{Weitere Bemerkungen}
{\index{NIBD|Z}\index{NIBP|Z}
Der \enquote{Nicht-invasive Blutdruck} wird auf manchen
Überwachungsmonitoren mit \T{NIBD} oder \T{NIBP}
(\Lang{engl.} Non-Invasive Blood
Pressure) abgekürzt. }{}{}{}{}
}



\section{Apparative Untersuchungen}

\subsection{Pulsoxymetrie}
\label{topic:pulsoxymetrie}

% \subsubsection{Allgemein}

\Text{\TxBeschreibung}
{% 
Die Pulsoxymetrie ist eine einfache apparative
Untersuchungsmethode zur nicht-invasiven Ermittlung der
Sauerstoffsättigung des kapillären Blutes, welche mit einem
\E{Pulsoxymeter} durchgeführt wird. Sie dient zur
nicht-invasiven Überwachung der \T{Sauerstoffsättigung}
(\T{Sp\ce{O2}}) des Hämoglobins\footnote{Das Hämoglobin ist
der für die Sauerstoffbindung verantwortliche Stoff in den
roten Blutkörperchen (\prettyref{topic:haemoglobin}).} im
kapillären Blut. Das Messprinzip basiert auf einer Messung
\Z{der Lichtabsorption bzw. der Lichtremission} bei
Durchleuchtung der Haut mittels eines Sensors bestehend aus
einer Lichtquelle und einem Fotosensor. Auf diese Weise wird
das Verhältnis der mit Sauerstoff beladenen roten
Blutkörperchen zu den nicht beladenen bestimmt.
Als weiteren Wert zeigt das Pulsoxymeter die
\EE{Pulsfrequenz} an.

Als \EE{Messstelle} eignen sich meist leicht zugängliche
Körperstellen (\E{Finger, Zeh, Ohrläppchen}) -- die
Herstellerangaben über die Eignung des Sensors für die
gewünschte Messstelle müssen jedoch beachtet werden.
\cite{LoefflerBiochemie7,ArbeitstechnikenAZ1,Soehnke200607}

}{}{}{}{}
\TextSlide{Gerätearten}
{% 
Zur Messung der Sauerstoffsättigung mit dem Prinzip der
Pulsoxymetrie gibt es verschiedene Gerätearten:

\begin{itemize}
  \item \EE{Fingerclip-Geräte}: Diese vereinen Messgerät und
  Sensor in einem Gerät.
  \item \EE{Handpulsoxymeter}: Diese bestehen aus einem
  Handmessgerät und einem per Kabel mit dem Messgerät verbundenen Sensor.
  \item \EE{Funktionsmodul eines Überwachungsmonitors}:
  Viele Multifunktionsmonitore verfügen über die Möglichkeit zur Pulsoxymetrie
  mithilfe eines Sensors.
\end{itemize}

}{
\begin{itemize}
  \item Fingerclip-Geräte
  \item Handpulsoxymeter
  \item Funktionsmodul eines Überwachungsmonitors
\end{itemize}
}{}{}{}
\TextSlide{Sensortypen}
{%
Neben den Fingerclip-Geräten sind praktisch alle anderen
Gerätetypen mit wechselbaren, kabelgebundenen Sensoren
versehen. Diese Sensoren können je nach Anwendungsfall
ausgetauscht werden.

Prinzipiell kann man zwei Sensorgrundarten unterscheiden:
\begin{itemize}
  \item \EE{Einmalsensoren} \ldots{} sind üblicherweise
  Klebesensoren für diverse meist exponierte Körperstellen
  (Finger, Zeh, Stirn, Nasenspitze) und oft nur in
  Abhängigkeit der Patientengruppe (Säuglinge, Kinder,
  Erwachsene) zu verwenden. Zu beachten ist neben den schon
  genannten Faktoren Messstelle und Patientengruppe auch
  noch das \E{Ablaufdatum} und die maximale Verwendungsdauer
  von Einmalsensoren, da der Hersteller nur bei richtigem
  Einsatz für ein brauchbares Ergebnis garantiert.
  \item \EE{Mehrwegsensoren} \ldots{} sind ebenso für
  verschiede Körperstellen und in Abhängigkeit der
  Patientengruppe (Säuglinge, Kinder, Erwachsene)
  erhältlich.
\end{itemize}
}{
\begin{itemize}
  \item Einmalsensoren
  \item Mehrwegsensoren
\end{itemize}
}{}{}{}

\Text{Funktionsprüfung}
{%
Vor der Verwendung ist das jeweilige Gerät einer
Sichtprüfung auf offensichtliche Schäden zu unterziehen
(Gehäuse, Kabel, Sensor). Nach dem Einschalten ist der
Selbsttest des Gerätes abzuwarten und der Batteriestand zu
prüfen. Anschließend ist das Gerät im Selbstversuch zur
Messung der eigenen Sauerstoffsättigung und der Pulsfrequenz
zu verwenden. Nur bei entsprechenden Messwerten darf das
Gerät zur Messung am Patienten verwendet werden.
}{}{}{}{}

\Text{Anwendung}
{%
Zur Anwendung am Patienten wird eine entsprechende
Messstelle gesucht. Dabei ist besonders auf Nagellack,
künstliche Nägel und Verschmutzung der Messstelle zu achten.
Bei starker Verschmutzung muss die Messstelle zuerst
entsprechend gereinigt werden (z.B. Alkoholtupfer bei
Fingern). Dann wird das Gerät eingeschaltet und der
Selbsttest abgewartet. Danach wird der Sensor an der
entsprechenden Stelle platziert (meist Mehrwegsensoren für
Erwachsene zur Verwendung an Fingern).
Nach etwas Wartezeit können die Messwerte abgelesen werden.
Es muss natürlich immer der klinische Zustand und die
Anamnese mit in Betracht gezogen werden.


\Fazit{Eine ausreichende Blutversorgung der Körperstelle,
an der der Sensor angebracht ist, ist Voraussetzung für die
Pulsoxymetrie.}

Das Pulsoxymeter eignet sich sehr gut für das Monitoring,
ersetzt aber niemals das wachsame Auge auf den Patienten!
}{}{}{}{}

\TextSlide{Ermittelte Werte}
{%
Die Sauerstoffsättigung wird in Prozent angegeben und wird
mit \T{Sp\ce{O2}} abgekürzt. Der Normalwert liegt zwischen
\EE{\VonBis{96}{100}\PC*}. Bei Patienten mit COPD findet man
jedoch oft sehr niedrige Werte (z.\,B.
90\PC*), ohne dass die Patienten Beschwerden haben.

\bigskip

}{
\begin{itemize}
  \item Normalwert: \VonBis{96}{100}\PC*
  \item COPD-Patient: auch Werte < 90\PC* möglich ohne
  Beschwerden
\end{itemize}
}{}{}{}


\Text{Beurteilung}
{%
Bei der Beurteilung des Sp\ce{O2} muss berücksichtigt
werden, \E{ob und wieviel \EE{Sauerstoff} der Patient
bereits erhält}. Eine Sauerstoffsättigung von 95\PC* bei
einem Patienten ohne zusätzliche Sauerstoffgabe ist nicht
beunruhigend, wohingegen derselbe Wert bei einer
\E{Sauerstoffgabe} von 15{\LPerMin*} ein Hinweis auf eine
schwere Störung der Atmung darstellt.

\Fazit{Die Beurteilung der Sauerstoffsättigung muss
berücksichtigen, ob und wieviel Sauerstoff der Patient
bereits erhält.} 
% Auch ein
% Normalwert bedeutet nicht, dass der Patient genug Sauerstoff hat, beziehungsweise, dass keine
% Sauerstoffgabe erfolgen soll. Die Sauerstoffgabe richtet
% sich immer nach dem klinischen
% Zustandsbild des Patienten und dessen Diagnose, und nur in
% Ausnahmefällen (z.\,B. Asthma bronchiale,
% \prettyref{topic:asthma-bronchiale}) nach dem
% Sp\ce{O2}-Wert. 
Der Sp\ce{O2}-Wert soll außerdem auch im
Verlauf beurteilt werden.
\cite{Soehnke200607}

\Fazit{Behandle den Patienten und nicht das Gerät!}
}{
\begin{itemize}
  \item Bedenke: Wieviel \ce{O2} bekommt der Patient
  bereits?
\end{itemize}
}{}{}{}




\TextSlide{\TxGefahren{} und Fehlerquellen}
{%
Die größte Gefahr ist die falsche Interpretation von
Messwerten. 

\begin{itemize}
  \item Manche Geräte sind \EE{nicht defibrillationssicher},
  d.\,h. es können während und nach der Defibrillation
  falsche Werte angezeigt werden, das Gerät darf während der
  Defibrillation nicht berührt werden. Bei Mehrwegsensoren
  besteht das Risiko der Keimübertragung (Kreuzinfektion,
  \FullPrettyRef{topic:kreuzinfektion}).
  \item Eine besondere Fehlerquelle ist die
  \EE{Kohlenmonoxid-(\ce{CO})-Vergiftung}.
  Kohlenmonoxid-beladenes Blut ist ebenso hellrot wie
  Sauerstoff-beladenes, daher wird der Patient nicht
  zyanotisch obwohl er eigentlich zu wenig Sauerstoff hat.
  Das Pulsoxymeter verwechselt daher folglich auch
  Kohlenmonoxid mit Sauerstoff und zeigt falsch hohe Werte
  an!
  \item Voraussetzung für eine korrekte Messung ist eine
  \EE{ausreichende Durchblutung} der entsprechenden
  Messstelle, z.\,B. durch Kälte, Schock oder auch eine
  Blutdruckmessung kann die Messung gestört oder unmöglich
  gemacht werden. Weitere Hindernisse sind z.\,B. kalte
  Finger (verringerte Durchblutung) oder \EE{Nagellack}
  (Licht kann den Nagellack nicht durchdringen, bzw. dessen
  Farbe wird verfälscht.) Mit Nagellack-Entferner kann der
  Nagel eventuell gereinigt werden. Alternativ kann der
  Sensor auch um 90\Deg gedreht werden.
  \item Von der Seite \EE{einstrahlendes Licht} kann ebenso
  das Ergebnis verfälschen, da es sich ja um eine lichtabhängige
  Messmethode handelt.
\end{itemize}

}{
\begin{itemize}
  \item \ce{CO}-Vergiftungen (falsch hohe Werte)
  \item Minderdurchblutung der Peripherie; Entweder durch
  Zentralisation beim Schock oder Unterkühlung.
  \item Nagellack
  \item Künstliche Fingernägel
  \item Lichteinstrahlung von der Seite
\end{itemize}
}{}{}{}



\Text{Hygienische Wiederaufbereitung}
{%
Die hygienische Wiederaufbereitung muss immer nach
Herstellerangaben erfolgen.
Einmalsensoren sind nach der Messung zu entsorgen.
Die meisten Geräte und Sensoren dürfen mittels
\E{Wischdesinfektion mit Flächendesinfektionsmittel}
aufbereitet werden. Die Geräte sind üblicherweise nicht
wasserdicht, es darf kein Desinfektionsmittel in das Gerät
gelangen.
}{}{}{}{}








% \subsubsection{Nellcor{\protect\texttrademark} NPB-40,
% Nellcor{\protect\texttrademark} NPB-40 Oximax,
% Nellcor{\protect\texttrademark} N-65 Oximax}
% 
% 
% Diese Geräte sind \E{nicht} defibrillationssicher. Sie dürfen während der
% Defibrillation am Patienten bleiben, doch während und kurz nach dem Schock können
% die Messwerte ungenau sein und das Gerät darf während des Schocks nicht berührt
% werden.
% 
% \TODO{GABS}{2010-06-14}{POI- Pulsoxy Nellcor: Symbol
% \enquote{Batterie} und \enquote{durchgestrichener Lautsprecher} einfügen.}{}
% Eine leer werdende Batterie wird durch ein Batterie-Symbol gekennzeichnet. Die
% leeren Batterien fachgerecht entsorgen und durch 4 neue Batterien der Type \enquote{AA}
% unter Beachtung der Polarität (siehe Symbole im Batteriefach auf der Rückseite
% des Gerätes) ersetzen. Die Lautstärke des Pulstones kann durch die Tasten \Sign{↑} und \Sign{↓}
% eingestellt werden.
% 
% Bei Auftreten eines Alarmes kann dieser für max. 120 Sekunden 
% \A{durch Drücken der Taste []}
% stummgeschalten werden.
% 
% 
% 
% 
% \subsubsection{Masimo{\protect\texttrademark} Rad5v}
% 
% Dieses Gerät darf während der Defibrillation verwendet werden, es kann jedoch während des
% Schocks und 20 Sekunden danach zu verfälschten Messwerten kommen.
% Die Batteriezustandsanzeige ist unterhalb der Ein-/Aus-Taste in Form von 4 grünen LEDs
% angebracht. Wenn nur mehr eine grüne LED der Batteriezustandsanzeige angezeigt wird, sind
% die leeren Batterien fachgerecht zu entsorgen und durch 4 neue Batterien der Type „AA“
% unter Beachtung der Polarität (siehe Symbole im Batteriefach auf der Rückseite des Gerätes)
% ersetzen.
% 
% \TODO{GABS}{2010-06-14}{Pulsoxy Masimo Rad5v: Symbol \enquote{Taste Pulston}
% einfügen.}{}
% \opt{STATUS-DRAFT}{\A{Draft:}
% Die Lautstärke des Pulstones kann durch Drücken der Taste [] in 4 Stufen eingestellt werden
% (Aus, Laut, Mittel, Leise).
% }
% Dieses Gerät alarmiert nicht bei niedriger Sauerstoffsättigung, lediglich bei fehlendem
% Messsignal und eignet sich somit nur bedingt für Überwachungsaufgaben.
% 
% 

\subsection{Blutzuckermessung}
\label{topic:blutzucker-messung}

\TextSlide{Interpretation und Normalwert}
{
Der Blutzuckerspiegel wird in \E{Milligramm pro Deziliter}
(\MgDl\Z{, entspricht \MetricUnit{mg\,\PC}})\Z{, manchmal 
auch in \E{Millimol pro Liter} (\MetricUnit{\nicefrac{mmol}{\Liter}}),} 
gemessen. Der
Wert ist abhängig von der letzten Mahlzeit des Patienten.
Nach der Nahrungsaufnahme steigt der Blutzuckerspiegel stark
an und fällt dann wieder auf einen \enquote{nüchternen
Normalwert} ab. Der normale \EE{Nüchtern-Blutzuckerspiegel}
bewegt sich zwischen \EE{\Range{80}{100}\MgDl*} 
\Z{(\Range{4,4}{5,6}\,\MetricUnit{\nicefrac{mmol}{\Liter}})}. Im
Akutfall kann der Normalbereich jedoch großzügiger
angesetzt werden, da bei Werten zwischen 60 und 200\MgDl*
keine Notfallsymptome auftreten.
\cite{HeroldInnereMedizin2012}
}{
\begin{itemize}
    \item Abhängig von letzter Mahlzeit!
    \item Normalwert: \Range{80}{100}\MgDl*
    (nüchtern)
    \item Bereich ohne Notfallsymptome:
    \Range{60}{200}\MgDl*
%     \item Notfallrelevanter Hypoglykämie: <
%     60\MgDl*
%     \item Notfallrelevante Hyperglykämie: >
%     200\MgDl*
\end{itemize}
}{}{}{}

\Text{Material}
{
\begin{enumerate}
  \item Handschuhe
  \item Kontamed-Box
  \item Alkoholtupfer
  \item Trockene Tupfer
  \item Blutzuckermessgerät
  \item Blutzuckermessstreifen
  \item Lanzette
\end{enumerate}
}
{}
% {./images/gabriel-sebastian-aass/bz-wert_1214-00800.jpg}
% {\enquote{90\MgDl*.}}
% {GABS}


\Text{Durchführung}
{
%\begin{multicols}{2}
\begin{enumerate}
  \item Finger auswählen: möglichst von der nicht dominanten
  Hand.
  \item \StepHeading{Einstichstelle} bestimmen: Am Fingerendglied,
  seitlich (Rand!), nicht in die Fingerkuppe, nicht zu nahe
  am Gelenk. Außenseite des Ringfinger bevorzugen.
  \item Stelle \StepHeading{desinfizieren}: Mit Alkoholtupfer abwischen
  und gemäß Einwirkzeit einwirken lassen.
  \item \StepHeading{Blutzuckermessstreifen halb in das Gerät} schieben,
  aber noch nicht einrasten lassen.
  \item Mit einer schnellen, beherzten Bewegung den
  Patientenfinger \StepHeading{mit der Lanzette stechen} und sofort
  herausziehen \E{(\enquote{Rein-raus-Methode})}. 
  Bei der Verwendung von Lanzettenautomaten (z.\,B. mit Federdrucksystemen) 
  sind die herstellerspezifischen Anwendungshinweise zu beachten.
  \item Lanzette in die \E{Kontamed-Box} \StepHeading{entsorgen}.
  \item Ersten Bluttropfen mit trockenem Tupfer \StepHeading{abwischen}.
  \item \StepHeading{Messstreifen} ganz in das Gerät schieben, das Gerät
  schaltet sich ein und führt einen Selbsttest durch. Wenn
  im Display ein Tropfen angezeigt wird ist das Gerät
  messbereit.
  \item Den Messstreifen mit dem Gerät \StepHeading{an den Bluttropfen halten}, 
  das Blut steigt aufgrund der Kapillarwirkung auf. 
  \hfill\Commentary[Blutzuckermessung, Kapillarwirkung]{Je 
  nach Modell kann es hier zu Abweichungen kommen, so gibt es z.\,B. 
  Geräte, bei denen der Teststreifen mit Blut benetzt werden muss. Im 
  professionellen Umfeld sollten jedoch aus hygienischen Gründen 
  nur Geräte verwendet werden, 
  welche nach dem Kapillar-Prinzip verfahren.}
  \item Beim \StepHeading{Piepton} ist genug Blut in die Testkammer
  eingeströmt und das Gerät beginnt mit der Messung und kann
  entfernt werden. 
  \item Mit einem trockenen \StepHeading{Tupfer} auf die Einstichstelle
  drücken (lassen).
  \item Ein neuerlicher Piepton zeigt an, dass die Messung
  beendet wurde und der Wert vom Display abgelesen
  werden kann. 
  \item \StepHeading{Messstreifen entsorgen} (normaler Restmüll).
\end{enumerate}
%\end{multicols}
}{

}{}{}{}

% \Picture
% 	{<Titel>}%1
% 	{<Beschreibung>}%2
% 	{<Label>}%3
% 	{<Quelle>}%4
% 	{<Lizenz>}%5
% 	{<Datei>}%6
% 	{<Breitenmodifizierer>}%7
% 	{<Float>}%8
\Picture
{Material zur Blutzuckermessung}%1
{}%2
{}%3
{Ch. Pallinger}%4
{MFG}%5
{./images/pallinger-christoph-aass/Blutzucker_32796_export_01312.jpg}%6
{}%7
{H}%8

\begin{PictureSeries}{}{Bilderserie: Blutzuckermessung}
\PictureSeriesRow{2}
{./images/motal-aass/00800/neuro6.jpg}
{Blutzuckermessung mit Messgerät mit Kapillar-Methode \SourcePicture{Motal}}
{./images/gabriel-sebastian-aass/bz-wert_1214-00800.jpg}
{Ergebnis: \enquote*{90\MgDl*.} \SourcePicture{GABS}}
{empty}
{}
{}
{}
\end{PictureSeries}






\subsection{Körpertemperatur}
\label{topic:koerpertemperatur-messung}

Die Körpertemperatur kann an unterschiedlichen Stellen und
mit unterschiedlichen Mitteln bestimmt werden. 

\TextSlide{Messarten}
{%
Mit dem
\EE{herkömmlichen Fieberthermometer}\footnote{Unter
herkömmlichen Fieberthermometern versteht man heute
elektronische Messgeräte. Quecksilberthermometer werden
wegen des darin enthaltenen, giftigen Schwermetalls
Quecksilber nicht mehr verwendet!} kann die Temperatur unter
der \E{Achsel} (\E{axillär}) oder im \E{Enddarm}
(\E{rektal}, vor allem bei Kleinkindern) bestimmt werden, 
die Messung im Mund ist im professionellen Umfeld 
eher unüblich. Es ist
bei jeder der Messung eine Schutzhülle zu verwenden.

Mit dem \EE{Ohrthermometer} kann mittels Infrarot die
Temperatur im \E{Mittelohr} bestimmt werden. Es sind die
entsprechenden \E{Schutzkappen} zu verwenden.
Die Temperatur im Enddarm und im Mittelohr bezeichnet man
als \T{Körperkerntemperatur}, diese beträgt um die
37\,\DegC. Die unter der Achsel gemessene Temperatur kann
ca. \nicefrac{1}{2}\DegC* niedriger sein.
}{
\begin{itemize}
  \item nach dem Ort der Messung
  \begin{itemize}
    \item axillär
    \item rektal
  \end{itemize}
  \item nach dem Gerätetyp
  \begin{itemize}
    \item herkömmliches (elektronisches) Thermometer
    \item Ohrthermometer
  \end{itemize}
\end{itemize}
}{}{}{}

\begin{PictureSeries}{}{Messung der Körpertemperatur}
\PictureSeriesRow{2}
{./images/ASBOE-Grassl-gjh-1/CRW_0901-00441pt}
{Temperaturmessung im Ohr eines Kindes \SourcePicture{ASBÖ/gjh}}
{./images/gabriel-sebastian-aass/fieberthermometer-00441pt.jpg}
{Fieber. Quecksilberthermometer dürfen nicht mehr verwendet und müssen fachgerecht entsorgt werden. \SourcePicture{Gabriel}}
{./images/ASBOE-Grassl-gjh-1/}
{}
{}
{}
\end{PictureSeries}



\subsection{EKG}
\label{topic:ekg}
\index{EKG}

Das \E{Elektrokardiogramm} (\Abk{EKG}) zeigt mittels eines
entsprechenden Gerätes die elektrische Herzaktivität des
Reizbildungs- und Reizleitungssystems an. Dazu werden auf
dem Körper nach einem bestimmten Muster Elektroden
aufgeklebt und mittels eines Kabels mit dem \Abk{EKG}-Gerät
verbunden. Im Rettungsdienst werden fast ausschließlich
Klebeelektroden für den einmaligen Gebrauch verwendet. 

\AuthNotes{GABS: Sollte im Defi-Teil durchgenommen
werden?:} 

\Layout{22}{\Z{EKG -- Ablauf einer Erregung}}
{\Z{\begin{itemize}
    \item P-Welle $\approx$ Vorhoferregung
    \item PQ-Zeit $\approx$ AV-Überleitung
    \item QRS-Komplex $\approx$ Kammererregung
    \item ST-Strecke + T-Welle $\approx$ Erregungsrückbildung
\end{itemize}
}
}
{}
{./images/wikipedia-cc-by-sa/EKG_Komplex}
{Eine elektrische Herzaktion}
{WMC:Hank, CC-BY-SA-2.0-DE}



\Cave{Das EKG zeigt nur die \E{elektrische} Aktivität,
nicht aber die tatsächliche Muskelarbeit an! 

\EE{Es erlaubt
keine Aussage über die Auswurfleistung! }}

\Fazit{Grundsatz: \Speech{Selbst ein pulsloser Patient kann
ein \enquote{schönes EKG haben}} (allerdings nicht
lange)!}

\subsubsection{EKG -- Ableitungen}
\label{topic:ekg-ableitungen}

\TextSlide{\TxBeschreibung}{
Grundsätzlich unterscheidet man zwischen den Extremitätenableitungen
und den Brustwandableitungen. 
\hfill\Commentary[EKG-Ableitungen]{Brustwandableitungen nach
\I{Einthoven} und \I{Goldberger}, Brustwandableitungen nach
\I{Wilson}. \ParBugfix
Neben diesen
Standardableitungen gibt es noch für spezielle Fragestellungen weitere
Ableitungmöglichkeiten wie z.\,B. die Ableitungen nach
\I{Nehb} oder die
\I[Brustwandableitungen!rechtsthorakale]{rechtsthorakalen} bzw.
\I[Brustwandableitungen!linksdorsale]{linksdorsalen
Brustwandableitungen}.
Zusammen ergeben die Elektroden der Frontal- und Horizontalebene 
eine dreidimensionale Darstellung der
elektrischen Herzaktivität: \\
  \EE{Inferior}: \EKG{II}, \EKG{III}, \EKG{aVF}\\
  \EE{Anterior / anteroseptal}: \VonBis{\EKG{V1}}{\EKG{V4}}\\
  \EE{Lateral}: \EKG{I}, \EKG{aVL}, \EKG{V5}, \EKG{V6}\\
  \EE{Dorsal}: \VonBis{\EKG{V7}}{\EKG{V9}}\\
  \EE{Rechtes Herz}: \VonBis{\EKG{V3R}}{\EKG{V6R}}\\
\cite{EkgIsabell5,InnereMedQuick3}}

\begin{itemize}
  \item Extremitätenableitungen\Z{: \EKG{I}, \EKG{II},
  \EKG{III}, \EKG{aVL}, \EKG{aVR}, \EKG{aVF}}
  \item Brustwandableitungen\Z{:
  \VonBis{\EKG{V1}}{\EKG{V6}}}
\end{itemize}
}{
\begin{itemize}
  \item Extremitätenableitungen
  \item Brustwandableitungen
\end{itemize}
}{}{}{}

Es haben sich folgende Begriffe eingebürgert:

\begin{description}
    \item[Extremitätenableitung] \enquote{\T{Rhythmusstreifen}}
    zur Beurteilung des Herzrhythmus und zum Monitoring
    \item[Extremitätenableitung\,+\,Brustwandableitung]
    (\enquote{\T{12-Kanal-EKG}}) -- Damit können Störungen
    der jeweiligen Herzregion zugeordnet werden. Wenn vorhanden,
    ist das 12-Kanal-\Abk{EKG} (\emph{\enquote{12er}}) für die
    Diagnostik das Mittel der Wahl. 
\end{description}



\TextSlide{Extremitätenableitungen}
{
Die \T{Extremitätenableitungen} werden mittels vier
Elektroden abgeleitet. Je eine Elektrode wird an den Enden der
Extremitäten angebracht. Zur Vereinfachung können die
Elektroden auch am Übergang vom Rumpf zu den Extremitäten
angebracht werden (z.\,B. statt dem rechten Arm an der
rechten Schulter). Wichtig ist, dass die Elektrodenpärchen
immer auf \EE{gleicher Höhe} angebracht werden!

Aus den Informationen der vier Elektroden erzeugt das
EKG-Gerät \EE{sechs Ableitungen} \Z{(\EKG{I}, \EKG{II},
\EKG{III}, \EKG{aVL}, \EKG{aVR} und \EKG{aVF}).}
}{
\begin{itemize}
  \item 4 Elektroden an den Extremitäten
  
  (bzw. am Übergang Rumpf -- Extremitäten).
  \item 6 Ableitungen.
\end{itemize}
}{}{}{}


\begin{table}[H]
\TableBaseModification
\Captionof{table}{Farbe und Positionen der EKG-Elektroden
für die Extremitätenableitung.}

\begin{tabular}{lll}
\addlinespace\toprule\addlinespace
& 
\TabHeading{Position} 
&
\TabHeading{Alternative Position}
\\\midrule
\colorbox{red}{\sffamily\bfseries\textcolor{white}{Rot}} & Rechte Hand & Rechte
Schulter
\\
\colorbox{yellow}{\sffamily\bfseries Gelb} & Linke Hand & Linke
Schulter
\\
\colorbox{green}{\sffamily\bfseries\textcolor{white}{Grün}} & Linker Fuß  &
Linke Leiste
\\
\colorbox{black}{\sffamily\bfseries\textcolor{white}{Schwarz}} & Rechter Fuß &
Rechte Leiste \\\bottomrule
\end{tabular}
\end{table}


\begin{figure}[H]
\begin{center}
\includegraphics[width=\linewidth]{./images/geraete/ECGcolor_edited}
\Captionof{figure}{\AuthNotesPicLegalOk{}{PD}Extremitätenableitungen
und Brustwandableitungen.
\SourcePicture{WmCo/Madhero88, PD}}
\end{center}
\end{figure}




\TextSlide{Brustwandableitungen}
{
Die \T{Brustwandableitungen} werden mittels sechs Elektroden
abgeleitet und erlauben eine genauere Zuordnung von
Störungen zu der jeweiligen Region des Herzens. Aus den
Informationen der sechs Elektroden erzeugt das
\Abk{EKG}-Gerät \EE{sechs Ableitungen}:
\Z{\VonBis{\EKG{V1}}{\EKG{V6}}}.
}{
\begin{itemize}
    \item 6 Elektroden am Brustkorb.
    \item 6 Ableitungen.
\end{itemize}
}{}{}{}





\begin{table}[H]
\TableBaseModification
\Z{
\Captionof{table}{Position der EKG-Elektroden bei der Brustwandableitung.}

\begin{tabulary}{\linewidth}[H]{CL}
\addlinespace\toprule\addlinespace
\noindent\TabHeading{Abl.}
 &
\noindent\TabHeading{Position}
\\\addlinespace\midrule\addlinespace
\EKG{V1} & 4. Zwischenrippenraum rechts vom Sternum
\\\addlinespace
\EKG{V2} & 4. Zwischenrippenraum links vom Sternum
\\\addlinespace
\EKG{V3} & Mitte zwischen \EKG{V2} und \EKG{V4}
\\\addlinespace
\EKG{V4} & 5. Zwischenrippenraum links in der Verlängerung
der Schlüsselbein-Mittellinie
\\\addlinespace
\EKG{V5} & 5. Zwischenrippenraum vordere Achsellinie (zw. \EKG{V4} und \EKG{V6}) 
\\\addlinespace
\EKG{V6} & 6. Zwischenrippenraum in der mittleren
Achsellinie 
\\\addlinespace\bottomrule
\end{tabulary}
}
\end{table}


Die richtige Positionierung der
\Abk{EKG}-Elektroden ist Voraussetzung um ein
sinnvoll auswertbares und den einzelnen
Herzabschnitten zuordenbares \Abk{EKG} zu erstellen.
Falsch angebrachte Elektroden können ein \Abk{EKG} derart
verfälschen, dass es zu gravierenden Fehldiagnosen kommen
kann. \hfill\Commentary[EKG, Positionierung der EKG-Elektroden]
{%
Zum Beispiel Rahmen einer Akut-Angiographie im Zuge eines
Myokardinfarktes sind die dokumentierten EKGs (und damit
auch die im Rettungsdienst abgeleiteten) ein wichtiges
Hilfsmittel, um das in der Akutsituation maßgebliche
Koronargefäß zu identifizieren. 
\ParBugfix
In \cite{Schnelle200607} wurde untersucht,
wie sich die Anbringungsstelle der Elektroden auf das
\Abk{EKG} auswirken kann.}

\bigskip

\begin{figure}[H]
\includegraphics[width=\linewidth]{./images/wikipedia-pd/12_lead_generated_sinus_rhythm}
\caption{Unauffälliger 12-Kanal-EKG-Befund \SourcePicture{WM/PD}}
\end{figure}



\begin{Literary}
  \fullcite{ElektroKomikoGraphie7}

  \fullcite{ekg-online.de}
  
  \fullcite{EKG-Knacker2}

  \fullcite{EkgIsabell5}
  
  Sowie:
  \cite{EcgEmergencyDec2,Harrisons18DeKap228,KuehnEkgFortbildung02T1,KuehnEkgFortbildung03T2}
  
\end{Literary}


\section{Körperliche Untersuchungen}

\subsection{Abtasten des Abdomens}

\TextSlide{Abtasten des Abdomens}
{\label{topic:untersuchung-abdomen-abtasten}
Das Abdomen kann man abtasten um die Anspannung der
Bauchdecke zu beurteilen, Schmerzen zu beurteilen und
zuordnen und um Fremdkörper zu ertasten.
Man beginnt damit, den Patienten zu fragen, ob und wo er
Schmerzen hat. Die Untersuchung beginnt man in jenem
Quadranten, welcher am weitesten von den Schmerzen entfernt
ist!

Man legt nun eine Hand flach auf den Bauch des Patienten und
die andere Hand flach auf die erste. Mit den Fingern drückt
man nun vorsichtig mehrere Zentimeter auf die Bauchdecke und
beurteilt, ob der Bauch weich oder verhärtet ist. Dabei
fragt man den Patienten ob er Schmerzen empfindet. Dies
wiederholt man in allen vier Quadranten.

Weiters kann man beurteilen ob der Patient Schmerzen durch
den Druck hat (Druckschmerz) oder ob er Schmerzen beim
Loslassen bekommt
(\T{Loslassschmerz}\index{Loslassschmerz}).
}
{
\begin{itemize}
  \item Fremdkörper? Wunden?
  \item Schmerzen?
  \item Einteilung in 4 Quadranten
  \item Abtasten der 4 Quadranten, beginnen dort wo kein
  Schmerz
  \item harter oder weicher Bauch?
  \item Druck- oder Loslassschmerz?
\end{itemize}
}{}{}{}

\begin{PictureSeries}{}{Inspektion und Abtasten des Abdomens \SourcePicture{ASBÖ/gjh}}
\PictureSeriesRow{2}
{./images/ASBOE-Grassl-gjh-1/IMG_0423-00441pt}
{Abtasten des Abdomens {\ldots}}
{./images/ASBOE-Grassl-gjh-1/IMG_0412-00441pt}
{{\ldots} in allen 4 Quadranten}
{}
{}
{}
{}
\end{PictureSeries}

% \clearpage % TEMP temp clearpage
\subsection{Neurocheck}
\label{topic:neurocheck}
\index{Neurocheck}
%\begin{multicols}{2}

\Text{{\TxBeschreibung}}
{\DefinitionInPara{Der Neurocheck ist eine einfache, überblickshafte (\enquote{grobe})
Untersuchung der wichtigsten Funktionen des Nervensystems.}
%
Teile des Neurochecks werden während des
Einschätzungsblockes
(\FullPrettyRef{topic:einschaetzungsblock})
durchgeführt.}




\TextSlide{Neurocheck}{
\begin{enumerate}
  \item \EE{Bewusstseinslage}: Für eine grobe Einschätzung
  ist folgende Einteilung ausreichend\footnote{Diese
  Einteilung entspricht weitgehend dem \T{WASB}-Einteilung
  (\EE{W}ach, Reaktion auf \EE{A}nsprechen, auf
  \EE{S}chmerzreiz, \EE{b}ewusstlos)}:
  \begin{enumerate}
    \item bewusstseinsklar
    \item bewusstseinsgetrübt (somnolent, soporös)
    \item bewusstlos
  \end{enumerate}
  Die Glasgow Coma Scale (GCS) erlaubt eine
  detailliertere Beschreibung des Bewusstseins und des
  Verlaufs.
  Vgl.
  \FullPrettyRef{topic:bewusstsein}, und
  \FullPrettyRef{topic:gcs}
  \item \EE{Orientierung}:
  \begin{enumerate}
    \item zur Person (Patient kennt seinen Namen)
    \item zum Ort (Patient weiß wo er sich befindet)
    \item zur Zeit (Patient kann das ungefähre Datum, die
    ungefähre Uhrzeit bzw. Jahres-/Tageszeit nennen)
    \item Situation
  \end{enumerate}
  \item \EE{Pupillen} Begutachtung und Lichtreaktion
  \begin{enumerate}
    \item Pupillen gleich weit?
    \item Blickstarre (\T{Herdblick})? \index{Herdblick}
    \item \T{Lichtreaktion:} In jedes Auge zwei mal leuchten
    und jeweils das gleiche und das andere Auge beobachten
    
    Beurteilung von:
    \begin{enumerate}
      \item Schnelligkeit der Reaktion: prompt, verzögert?
      \item Reagieren immer beide
      Pupillen?~\Z{\footnote{\Z{isokor}}}
    \end{enumerate}
  \end{enumerate}
  \item Lichtscheu und Sehstörungen?
  \item \EE{Kraftprobe} grobe Prüfung an oberer und
  unterer Extremität, beurteilt wird die Kraft und die
  Seitengleichheit.
  \begin{enumerate}
    \item Patient überkreuzt die Hände geben, bei kräftigen
    Patienten nur Zeige- und Mittelfinger.
    \item Patienten bitten kräftig zuzudrücken.
    \item Mit Handflächen von unten gegen die Zehen drücken
    und Patienten bitten \emph{\enquote{auf das Gaspedal zu
    drücken}}.
    \item Von oben auf die Zehen drücken und Patienten bitten
    die Zehen Richtung Kopf zu heben.
  \end{enumerate}
  \item
  \T[Nackensteifigkeit!Untersuchung]{Nackensteifigkeit}\label{topic:nackensteifigkeit-untersuchung}:
  Die Nackensteifogkeit ist ein
  \I[Meningismuszeichen!Nackensteifigkeit!Untersuchung]{Meningismuszeichen}.
  Der Patient liegt \EE{flach} auf dem Rücken. Der
  Untersucher fasst den Kopf des Patienten und versucht das
  Kinn an die Brust zu führen. Es handelt sich um eine
  \EE{passive Bewegung}, der Patient darf dabei nicht
  mithelfen.
  Normalerweise ist diese Bewegung ohne Probleme möglich.
  
  Wenn eine Nackensteifigkeit vorliegt ist die Beugung nicht
  möglich und der Versuch schmerzhaft, der Patient bewegt
  evtl. den ganzen Rumpf mit. Die Bedeutung der
  Nackensteifigkeit wird unter \prettyref{topic:meningitis}
  besprochen. \EE{Bei jedem Patienten mit Fieber ist auf das
  Vorliegen einer Nackensteifigkeit zu untersuchen}, jedoch
  \Ca{nicht bei Verdacht auf Wirbelsäulenverletzung!}
  \item
  \EE{Blutzuckermessung}: Eine Blutzuckermessung gehört zu
  einem vollständigen Neurocheck.
\end{enumerate}
}{
\begin{itemize}
  \item Bewusstseinslage
  \item Orientierung
  \item Pupillen
  \item Lichtscheue, Sehstörungen
  \item Kraftprobe
  \item Nackensteifigkeit
  \item[+] Blutzuckermessung
\end{itemize}

Teile des Neurochecks werden während des
Einschätzungsblockes
durchgeführt.
}{}{}{}

%\end{multicols}


\begin{PictureSeriesWide}{}{Bilderserie: \EE{Neurocheck}
\SourcePicture{Motal}}
\PictureSeriesRowWide{3}
{./images/motal-aass/00800/neuro1.jpg}
{Pupillenreaktion, Lichtscheue, Sehstörungen }
{./images/motal-aass/00800/neuro2.jpg}
{Kraftprobe an den Armen}
{./images/motal-aass/00800/neuro3b.jpg}
{Kraftprobe an den Beinen}
{}
{}
\PictureSeriesRowWide{3}
{./images/motal-aass/00800/neuro4.jpg}
{Nackensteifigkeit}
{./images/motal-aass/00800/neuro5.jpg}
{Nackensteifigkeit}
{./images/motal-aass/00800/neuro6.jpg}
{Blutzuckermessung}
{}
{}
\end{PictureSeriesWide}


% FEHLT: Bild für Orientierung (stattdessen jetzt 2x
% Nackensteifigkeit)

\clearpage % TEMP temp clearpage
\subsection{Traumacheck}
\label{topic:traumacheck}%

\Text{{\TxBeschreibung}}
{\DefinitionInPara{Ein \T{Traumacheck} ist eine überblickshafte 
Suche nach Verletzungen oder Verletzungsanzeichen am (liegenden)
Patienten.}
%
Teile des Traumachecks werden bereits während des
Einschätzungsblocks
(\FullPrettyRef{topic:einschaetzungsblock}) im
Rahmen der \II[Schnelle
Trauma-Untersuchung]{Schnellen-Trauma-Untersuchung}
(\I{STU}) durchgeführt, um vital bedrohliche 
Verletzungen rechtzeitig zu erkennen.  
Bei vital bedrohten (zeitkritischen) Patienten wird
normalerweise \E{kein vollständiger} Traumacheck 
durchgeführt. 
}


\TextSlide{\E{Vollständiger} Traumacheck}
{%
\begin{itemize}
  \item Patient orientiert? Unfallhergang erinnerlich? Schmerzen?
  \item Kontrolle von Mund, Nase und Ohren auf Blutungen (mit einer
  Lampe)
  \item Pupillenkontrolle (mit einer
  Lampe)
  \item Inspektion und Abtasten mit festem Griff von Kopf bis Fuß. 
  \EE{Schmerzhafte Stellen zuletzt
  untersuchen.}
  
  Zur Inspektion muss der Patient \EE{entkleidet} werden, 
  bei immobilisierten oder bewusstlosen Patienten muss 
  dazu oft die Kleidung aufgeschnitten werden. 
  \Ca{Wärmeerhalt beachten!} Ggfs. den Patienten erst in 
  geschützter Umgebung (Fahrzeug, Behandlungsraum, \ldots) 
  entkleiden und Inspektion nachholen. 
  \begin{itemize}
    \item Schädel: Schädeldach, Hinterkopf,
    Stirn, Wangenknochen, Nasenbein, Kiefer
    \item Hals und Halswirbelsäule 
    (Tastschmerz?) 
    
    \Ca{Keine Manipulation bei Verdacht auf Wirbelsäulenverletzung!}
    \item Obere Extremität: Schultern, Arme, Hände
    \item Brustkorb: vorne und seitlich 
    (Prellmarken?
    Atemmechanik? 
    Hautemphysem? 
    Gegen leichten Druck einatmen lassen)
    \item Bauch in 4 Quadranten (Prellmarken?
    Abwehrspannung/harte Bauchdecke?),
    \item Becken von der Seite andrücken, schmerzhaft?
    \item Rücken und Wirbelsäule (Tast-, Klopfschmerz?
    
    \Ca{Keine Manipulation bei Verdacht auf Wirbelsäulenverletzung!}
    \item Untere Extremität: Beine (sichtbare Verletzungen?), Füße:
    Schuhe und Socken ausziehen)
  \end{itemize}
  \item Durchblutung, Motorik, Sensibilität (\T{DMS})\label{topic:dms}
  \begin{itemize}
    \item \label{topic:dms} Finger und Zehen: Kontrolle von
    \E{Durchblutung}
    \E{Motorik} (bewegen lassen) und 
    \E{Sensibilität}
    (\Speech{\enquote{Spüren Sie das?}}) 
    \item \T{Nagelbettprobe} (\EE{Rekapillarisierungszeit}):
    leichter Druck auf Nagelende bis Nagelbett weiß ist,
    loslassen. Zeit bis Nagelbett wieder rosarot ist sollte
    nicht länger als \VonBis{1}{2} Sekunden sein.
    Seitenvergleich! Eine verlängerte Rekapzeit deutet auf
    eine Durchblutungsstörung hin. (einseitig oder
    beidseitig? lokal oder systemisch?)
  \end{itemize}
  \item
  \T{Tupfertest}\index{Tupfertest}\label{topic:tupfertest}:
  etwas Blut aus Ohren oder Nase mit einem Tupfer auffangen.
  Sollte Liquor im Blut enthalten sein bildet sich um den
  roten Blutfleck ein heller, bernsteinfarbener Hof.
\end{itemize}
Bei Hinweis auf eine Schädelverletzung ist ein
\E{Neurocheck} (\FullPrettyRef{topic:neurocheck}) auf jeden
Fall erforderlich! }{
\begin{itemize}
  \item Bewusstsein, Orientierung
  \item Mund, Nase, Ohren
  \item Pupillen
  \item Inspektion und Abtasten: Von Kopf bis Fuß
  \begin{itemize}
    \item Schmerzen
    \item Blutungen
    \item Fehlstellungen
    \item DMS
    \item Atemmechanik
    \item Harte Bauchdecke
  \end{itemize} 
  \item Tupfertest
\end{itemize}

Teile des Traumachecks werden während des
Einschätzungsblocks
bei der \E{Schnellen-Trauma-Untersuchung}
(\E{STU}) durchgeführt.  

Bei vital bedrohten (zeitkritischen) Patienten wird
normalerweise \E{kein vollständiger} Traumacheck 
durchgeführt! 
}{}{}{}

\begin{PictureSeriesWide}{}{Die wichtigsten Schritte der \EE{Schnellen Trauma-Untersuchung} (\EE{STU}) \SourcePicture{ASBÖ/gjh}}
\PictureSeriesRowWide{3}
{./images/ASBOE-Grassl-gjh-1/Uni_14-00441pt}
{Zur Inspektion muss der Patient entkleidet werden. 
Bei immobilisierten oder bewusstlosen Patienten muss 
dazu oft die Kleidung aufgeschnitten werden. 
\Ca{Wärmeerhalt beachten!}}
{./images/ASBOE-Grassl-gjh-1/IMG_0310-00441pt}
{Begutachtung der Mundhöhle (hier anfangs unter 
manueller HWS-Fixierung)}
{./images/ASBOE-Grassl-gjh-1/IMG_0388-00441pt}
{Bei Verdacht auf eine Wirbelsäulenverletzung muss eine 
entsprechende Fixierung der HWS auch während der Untersuchung 
erfolgen (hier nicht weiter dargestellt).}
{}
{}
\PictureSeriesRowWide{3}
{./images/ASBOE-Grassl-gjh-1/IMG_0380-00441pt}
{Thoraxkompression seitlich {\ldots}}
{./images/ASBOE-Grassl-gjh-1/IMG_0381-00441pt}
{{..} und frontal}
{./images/ASBOE-Grassl-gjh-1/IMG_0423-00441pt}
{Abtasten des Abdomens {\ldots}}
{}
{}
\PictureSeriesRowWide{3}
{./images/ASBOE-Grassl-gjh-1/IMG_0412-00441pt}
{{\ldots} in allen 4 Quadranten}
{./images/ASBOE-Grassl-gjh-1/IMG_0420-00441pt}
{Beckenkompression}
{./images/ASBOE-Grassl-gjh-1/IMG_0426-00441pt}
{Abtasten des Oberschenkels}
{}
{}
\end{PictureSeriesWide}





% \begin{PictureSeriesWide}{}{Bilderserie: \EE{Traumacheck}.
% Achte auf den Wärmeerhalt! In kühler Umgebung soll ein
% ausführlicher Traumacheck erst im (geheizten) Fahrzeug
% durchgeführt werden.
% \SourcePicture{Motal}}
% \PictureSeriesRowWide{3}
% {./images/motal-aass/00800/traumacheck1.jpg}
% {Kleidung entfernen }
% {./images/motal-aass/00800/traumacheck2.jpg}
% {Schuhe ausziehen (Fixierung, Stiefelgriff) }
% {./images/motal-aass/00800/traumacheck3b.jpg}
% {Untersuchung des Kopfes }
% {}
% {}
% \PictureSeriesRowWide{3}
% {./images/motal-aass/00800/traumacheck4.jpg}
% {In alle Körperöffnungen im Kopf leuchten, Pupillenkontrolle.}
% {./images/motal-aass/00800/traumacheck5b.jpg}
% {Neben Blut kann auch Liquor austreten (Tupfertest).}
% {./images/motal-aass/00800/traumacheck6b.jpg}
% {Untersuchung des Thorax (inkl. Atemmechanik)}
% {}
% {}
% \PictureSeriesRowWide{3}
% {./images/motal-aass/00800/traumacheck7b.jpg}
% {Untersuchung des Abdomen und des Beckens}
% {./images/motal-aass/00800/traumacheck8b.jpg}
% {Untersuchung der Extremitäten (DMS, Nagelbettprobe) }
% {./images/motal-aass/00800/traumacheck9b.jpg}
% {Nach dem Traumacheck den entkleideten Patient vor
% Wärmeverlust schützen. }
% {}
% {}
% \end{PictureSeriesWide}
% 
% 
% 







% \section*{\protect\dsliterary{} Weiterführende Literatur}
\begin{Literary}
\fullcite{DahmerAnamneseBefund10}

\fullcite{LehmeyerAnamUnter1}

\fullcite{ChecklisteAnamUnt1}

\fullcite{BatesUntersuchung1}

\fullcite{FamPropPku2012}
\end{Literary}
% \begin{WidePar}
% \begin{multicols}{3}
% \small
% \begin{description}%{mmm}
%   \item[\cite{DahmerAnamneseBefund10}] \emph{\bibentry{DahmerAnamneseBefund10}}
%   \item[\cite{LehmeyerAnamUnter1}] \emph{\bibentry{LehmeyerAnamUnter1}}
%   \item[\cite{ChecklisteAnamUnt1}] \emph{\bibentry{ChecklisteAnamUnt1}}
%   \item[\cite{BatesUntersuchung1}] \emph{\bibentry{BatesUntersuchung1}}
%   \item[\cite{FamPropPku2011}] \emph{\bibentry{FamPropPku2011}}
% \end{description}
% \end{multicols}
% \end{WidePar}
% 



\ChapterEnd


